\section{Control estructural, video y la próxima generación de arquitecturas generativas}

La generación de imágenes no se limita al condicionamiento por texto. Las señales estructurales como Pose, edge y scribble tienen correspondencia espacial, mientras que en subject-driven generation y en editing puede no haber una alineación directa; el video añade además la dimensión temporal. Las últimas diapositivas amplían DiT de un backbone para imágenes individuales a una plataforma más general de conditional generation.

\subsection{El control estructural no es un problema único}

Si la señal de control corresponde a las posiciones de los píxeles de salida, puede usarse un auxiliary encoder que produzca features multiescala y las inyecte en el base model; si la condición es la identidad del sujeto o una instrucción de edición, la relación de posiciones no es fija y se requieren token interaction, adapter o mecanismos basados en recuperación. Esta sección establece una clasificación a partir de «si la condición tiene correspondence espacial», porque eso determina si la condición puede inyectarse posición por posición y también si la evaluación debe centrarse en el seguimiento del contorno, la preservación de la identidad o la semántica de la edición.

\lecturefigure{slide-59.jpg}{Ejemplos de control estructural: edge, depth y pose tienen correspondencia espacial con la imagen objetivo}{59}

Estas señales pueden verse como observaciones adicionales alineadas con el noisy latent. El modelo debe conservar el contenido del sujeto y, al mismo tiempo, seguir con rigor el contorno o la postura. La evaluación no puede limitarse a la estética; también deben calcularse control accuracy y content preservation.

\lecturefigure{slide-60.jpg}{ControlNet/PixArt-$\delta$: una red auxiliar codifica la condición visual y la inyecta en el base DiT}{60}

La auxiliary network copia o adapta parte del backbone e inyecta las features estructurales mediante un zero-initialized residual. Así puede congelarse el base model grande y reducirse el costo del ajuste fino. Para las subject/edit conditions sin correspondencia espacial, sumar directamente posición por posición no es razonable.

\teachervoice{01:08:47--01:10:40, el docente distingue en particular entre el control estructural y subject-driven/editing: el primero siempre tiene correspondence espacial y el segundo no necesariamente, por lo que la forma de inyección de ControlNet no puede generalizarse de manera mecánica.}

\begin{warningbox}{Conflicto entre la intensidad del control y la libertad de generación}
Una condición estructural demasiado fuerte fija la textura y la composición; una demasiado débil no logra que se siga el control. Deben reportarse las curvas fidelity--diversity con distintos valores de control scale, y revisarse el ruido en la condición y las posturas fuera del dominio.
\end{warningbox}

\subsection{Video: de tokens bidimensionales a interacción espaciotemporal tridimensional}

El video no solo tiene más cuadros: exige mantener la coherencia de identidad, movimiento y física a lo largo del tiempo. Aplicar joint attention a $T$ cuadros con $N$ tokens cada uno tiene un costo de $O((TN)^2)$; factorized attention es más económica, pero puede perder el acoplamiento espaciotemporal de largo alcance. Esta sección extiende la token accounting del DiT para imágenes individuales al eje temporal y compara qué problema resuelve cada una de estas opciones: full 3D interaction, la descomposición espacio--tiempo y la codificación posicional.

\lecturefigure{slide-61.jpg}{DiT para video: RoPE, 3D attention, factorized attention y retos de eficiencia}{61}

RoPE (Rotary Position Embedding) codifica la posición relativa rotando query/key, y puede extenderse a los ejes temporal y espaciales. Full 3D attention suele dar mejor calidad, pero con un costo altísimo; la atención descompuesta en espacio y tiempo es más controlable. Las métricas de video también deben cubrir la temporal consistency, no solo la nitidez de cada cuadro.

\teachervoice{00:59:10--01:01:38 y 01:10:40--01:11:44, en la sesión de preguntas se enfatiza que la dificultad del video no es «generar unas cuantas imágenes más, una por cuadro», sino el costo de la atención tridimensional y la coherencia temporal; el docente recomienda seguir líneas orientadas a la eficiencia como LTX-Video.}

\begin{knowledgebox}{Primera aparición: RoPE}
RoPE codifica la posición como una rotación de query/key, de modo que el producto punto incorpora de forma natural la información de posición relativa. En video pueden asignarse frecuencias de rotación al tiempo, la altura y el ancho, pero la extrapolación de longitud y la proporción entre los tres ejes todavía requieren diseño.
\end{knowledgebox}

\subsection{In-context generation y direcciones abiertas}

La próxima generación de modelos generativos busca leer ejemplos en contexto, como los modelos de lenguaje: dada una imagen de referencia, un par antes/después de una edición o una secuencia multimodal, el modelo infiere directamente la tarea. Los principales obstáculos de las arquitecturas actuales son secuencias demasiado largas, tipos de condición heterogéneos y la dificultad de construir datos de entrenamiento. Esta sección pasa de la conditional generation con interfaz fija a tareas guiadas por ejemplos; al leer las diapositivas siguientes hay que distinguir entre las demostraciones de un «formato de entrada unificado» y la evidencia de una «verdadera generalización entre tareas».

\lecturefigure{slide-62.jpg}{El problema de la próxima generación: cómo lograr in-context learning en modelos de difusión}{62}

In-context generation requiere colocar los ejemplos y el objetivo en un mismo contexto en el que puedan interactuar, sin ajuste fino para cada tarea. Si todas las imágenes se despliegan como tokens, el cómputo crece de forma explosiva; si se comprimen en exceso, ya no es posible copiar detalles o relaciones. La arquitectura debe combinar atención eficiente, representaciones jerárquicas y tokens de tarea.

\lecturefigure{slide-63.jpg}{Playground v3, FuseDiT, OmniGen y otros exploran la generación en contexto unificada}{63}

Estos modelos intentan organizar text, reference images y target latents en una secuencia unificada o en una interacción por etapas. Muestran una dirección, no una solución ya alcanzada al in-context learning general. Al compararlos, deben revisarse las tareas que admiten, la longitud de contexto, los datos de entrenamiento y si necesitan un adapter específico.

\lecturefigure{slide-64.jpg}{Diffusers ofrece implementaciones de varias arquitecturas, lo que facilita la reproducción y el hacking}{64}

La diapositiva de implementación destaca un punto de entrada a la investigación: una biblioteca unificada permite comparar módulos como DiT, PixArt, SANA y MMDiT, y reduce las diferencias ocultas de implementación. Sin embargo, «el código se ejecuta» no equivale a que el experimento sea reproducible; también se necesitan los pesos del modelo, el procesamiento de datos, la configuración de entrenamiento, las semillas aleatorias y los scripts de evaluación.

\lecturefigure{slide-65.jpg}{Conclusiones y direcciones no cubiertas: MoE, entrenamiento, post-training, alignment y evaluación}{65}

La última diapositiva enumera de forma explícita los vacíos: la mezcla de expertos (MoE) ya llegó a la generación de imágenes; el entrenamiento en sí es otro sistema completo; el posentrenamiento, la alineación de preferencias y la evaluación siguen siendo problemas independientes y difíciles. Este límite evita que confundamos el diseño del backbone con una inteligencia generativa completa.

\teachervoice{01:11:44--01:14:23, el docente considera la in-context generation un problema de la próxima generación de arquitecturas y reconoce explícitamente que no cubrió MoE, entrenamiento, post-training, preference alignment ni evaluation.}

\begin{importantbox}{La meta de la investigación en arquitecturas no es un sample más bonito}
Un sistema generativo completo también debe responder: ¿los datos tienen autorización y están deduplicados? ¿El entrenamiento es estable? ¿Cómo se alinea el modelo con las preferencias? ¿Las métricas automáticas son confiables? ¿El video tiene coherencia temporal? ¿El despliegue cumple con el presupuesto de latencia y de memoria de GPU? El backbone es solo una de esas capas.
\end{importantbox}

\subsection{Resumen de la sección}

El control estructural depende de si la condición está alineada espacialmente; el video introduce tokens tridimensionales y la coherencia temporal; la in-context generation exige que el modelo procese de manera unificada los ejemplos y el objetivo. Diffusers reduce la barrera de implementación, pero el entrenamiento, el post-training, MoE y la evaluación siguen siendo problemas de sistema abiertos.
